\label{prob:poisson-gamma}
Let $x_1, x_2, \ldots, x_N$ be an independent sample from $\Poiss(\lambda)$.
The prior distribution is the gamma distribution
\[
  p(\lambda) = \Gam(\lambda \mid a, b)
  = \frac{b^a}{\Gamma(a)} \lambda^{a - 1} \exp(-b \lambda),
  \quad a, b > 0.
\]
Find the Bayesian estimate of $\lambda$ as the mean of the posterior
distribution $p(\lambda \mid x_1, \ldots, x_N)$.

\begin{solution}
\solpart{Idea}
The numerator and the denominator of Bayes' theorem contain the same
function:
\[
  p(\lambda \mid X)
  = \frac{p(X \mid \lambda) \, p(\lambda)}
         {\int p(X \mid \lambda) \, p(\lambda) \, d\lambda},
  \qquad X = (x_1, \ldots, x_N).
\]
Do not calculate the integral directly:
\begin{align*}
  p(X \mid \lambda) \, p(\lambda) &= C \cdot g(\lambda)
    && \text{($C$ does not depend on $\lambda$)} \\
  g(\lambda) &= Z \, q(\lambda)
    && \text{($q$ is a known density)} \\
  \int C Z \, q(\lambda) \, d\lambda &= C Z
    && \text{(a density integrates to 1)} \\
  p(\lambda \mid X) &= \frac{C Z \, q(\lambda)}{C Z} = q(\lambda).
\end{align*}
The same method gives the posterior mean.

\solpart{Solution}
Let $S = \sum_{i=1}^{N} x_i$, $a' = a + S$, and $b' = b + N$.

\solstep{Step 1. Write the goal.}
\[
  \hat{\lambda}_{\mathrm{B}} = \mathbb{E}[\lambda \mid X]
  = \int \lambda \, p(\lambda \mid X) \, d\lambda.
\]

\solstep{Step 2. Find $p(\lambda \mid X)$ for Step 1.}
\[
  p(\lambda \mid X)
  = \frac{p(X \mid \lambda) \, p(\lambda)}
         {\int p(X \mid \lambda) \, p(\lambda) \, d\lambda}.
\]

\solstep{Step 2a. Separate the constant in the numerator of Step 2.}
\begin{align*}
  p(X \mid \lambda) \, p(\lambda)
  &= \prod_{i=1}^{N} e^{-\lambda} \frac{\lambda^{x_i}}{x_i!}
     \cdot \frac{b^a}{\Gamma(a)} \lambda^{a-1} e^{-b\lambda}
   = \frac{e^{-N\lambda} \lambda^{S}}{\prod_{i} x_i!}
     \cdot \frac{b^a}{\Gamma(a)} \lambda^{a-1} e^{-b\lambda} \\
  &= \underbrace{\frac{b^a}{\Gamma(a) \prod_{i} x_i!}}_{C}
     \cdot
     \underbrace{\lambda^{a' - 1} e^{-b'\lambda}}_{g(\lambda)}.
\end{align*}

\solstep{Step 2b. Find the known density for $g$ from Step 2a.}
\[
  \Gam(\lambda \mid a', b')
  = \frac{b'^{a'}}{\Gamma(a')} \lambda^{a' - 1} e^{-b'\lambda}
  \;\Longrightarrow\;
  g(\lambda) = \frac{\Gamma(a')}{b'^{a'}} \, \Gam(\lambda \mid a', b').
\]

\solstep{Step 2c. Substitute Steps 2a and 2b into Step 2.}
\[
  p(\lambda \mid X)
  = \frac{C \frac{\Gamma(a')}{b'^{a'}} \, \Gam(\lambda \mid a', b')}
         {C \frac{\Gamma(a')}{b'^{a'}}
          \int \Gam(\lambda \mid a', b') \, d\lambda}
  = \frac{C \frac{\Gamma(a')}{b'^{a'}} \, \Gam(\lambda \mid a', b')}
         {C \frac{\Gamma(a')}{b'^{a'}}}
  = \Gam(\lambda \mid a + S, \, b + N).
\]

\solstep{Step 3. Substitute Step 2c into Step 1.}
\[
  \mathbb{E}[\lambda \mid X]
  = \int \lambda \, \Gam(\lambda \mid a', b') \, d\lambda
  = \frac{b'^{a'}}{\Gamma(a')} \int \lambda^{a'} e^{-b'\lambda} \, d\lambda.
\]

\solstep{Step 3a. Find the integral for Step 3 with the same method.}
\[
  \lambda^{a'} e^{-b'\lambda}
  = \frac{\Gamma(a' + 1)}{b'^{a' + 1}} \, \Gam(\lambda \mid a' + 1, b')
  \;\Longrightarrow\;
  \int \lambda^{a'} e^{-b'\lambda} \, d\lambda
  = \frac{\Gamma(a' + 1)}{b'^{a' + 1}}.
\]

\solstep{Step 3b. Substitute Step 3a into Step 3.}
\begin{align*}
  \mathbb{E}[\lambda \mid X]
  &= \frac{b'^{a'}}{\Gamma(a')} \cdot \frac{\Gamma(a' + 1)}{b'^{a' + 1}}
   = \frac{a' \, \Gamma(a')}{\Gamma(a') \, b'}
     && \text{($\Gamma(a' + 1) = a' \, \Gamma(a')$)} \\
  &= \frac{a'}{b'} = \frac{a + S}{b + N}.
\end{align*}
\[
  \hat{\lambda}_{\mathrm{B}} = \frac{a + S}{b + N}.
\]

\solpart{Conclusion}
\begin{itemize}
  \item \emph{Conjugate prior:} gamma prior $\times$ Poisson likelihood
    $\to$ gamma posterior. The update only changes the parameters:
    $a \to a + S$ (events), $b \to b + N$ (observations).
  \item The estimate is a weighted mean of the prior mean and the
    estimate $\bar{x}$ from Problem~\ref{prob:poisson-mle}:
    \[
      \frac{a + S}{b + N}
      = \frac{a}{b + N} + \frac{S}{b + N}
      = \frac{b}{b + N} \cdot \frac{a}{b}
      + \frac{N}{b + N} \cdot \bar{x}
      \;\xrightarrow{N \to \infty}\; \bar{x}.
    \]
  \item If $a, b \to 0$, then $p(\lambda) \propto \lambda^{a - 1}
    e^{-b\lambda} \to 1/\lambda$, the non-informative prior from the notes
    below, and $\hat{\lambda}_{\mathrm{B}} \to S/N = \bar{x}$.
\end{itemize}
\end{solution}

\paragraph{Notes on comparing the estimates from
  Problems~\ref{prob:poisson-mle} and~\ref{prob:poisson-gamma}.}
\begin{enumerate}[label=\arabic*.]
  \item For more about non-informative priors, see Section~5.4 ``Priors''
    in Kevin P.\ Murphy, \emph{Machine Learning: A Probabilistic
    Perspective}.
  \item The prior $p(\lambda) \propto \frac{1}{\lambda}$ is a uniform prior
    on $\log \lambda$. To see this, use the change-of-variables rule for
    the density of a function of a random variable. See, for example,
    \href{https://en.wikipedia.org/wiki/Probability_density_function#Function_of_random_variables_and_change_of_variables_in_the_probability_density_function}{this
    article}.
\end{enumerate}
